\subsection{The Declaration Graph}
\label{sec:theorem-premise}

We now descend from modules to individual declarations. The module graph $G_{\mathrm{module}}$ of \S\ref{sec:import-graph} records which modules depend on which, treating each source file as an atomic unit; it cannot distinguish a module that contains one foundational definition from one that contains fifty routine lemmas. The \emph{declaration dependency graph} defined in this section resolves this limitation by recording the dependency of each individual declaration on its premises. At this finer granularity, the nodes are not files but theorems, definitions, and type constructors---the elementary units of mathematical content. Where the module graph captures the \emph{architecture} of the library, the declaration dependency graph captures the \emph{reasoning structure} of mathematics itself. In the next section (\S\ref{sec:namespace-graph}), we derive an intermediate family of graphs by aggregating declarations into namespaces; the declaration-level graph analyzed here is the ground truth from which those aggregations are computed.

The main findings of this section are threefold. First, the hub structure of the theorem--premise graph separates into two strata---a \emph{language infrastructure} layer (equality reasoning, numeral coercions) determined by the design of Lean's type system, and a \emph{mathematical infrastructure} layer (categories, real numbers, modules) determined by the logical structure of mathematics itself (\S\ref{sec:thm-types}--\S\ref{sec:thm-centrality}). Second, Louvain community detection recovers recognizable mathematical disciplines with moderate but imperfect alignment to the namespace hierarchy (NMI $= 0.34$), and only category theory achieves near-complete community--namespace correspondence (\S\ref{sec:thm-community}). Third, the cross-namespace density at the theorem level (50.9\%) far exceeds the module-level figure (37.1\%), revealing that mathematical reasoning violates cognitive classification more deeply than file organization alone suggests (\S\ref{sec:thm-cross-namespace}).

\paragraph{Data and Definitions}
\label{sec:thm-data}

The theorem--premise graph is constructed from two complementary data sources. The \texttt{lean4export} tool~\cite{lean4export} extracts the full environment of a Lean~4 project---every declaration, its type, its proof term, and the constants appearing in each---into a machine-readable format. The \texttt{lean-training-data} tool~\cite{lean_training_data} extracts premise-level dependencies by analyzing which constants appear in the proof term of each declaration. We combine these extractions into a single dataset, published on HuggingFace as \texttt{MathNetwork/MathlibGraph}.\xinze{LeanScout pipeline: could cross-validate}

\paragraph{What is a premise?}
When a mathematician writes a proof in Lean, the tactic script is compiled into a \emph{proof term}---a fully explicit expression in Lean's dependent type theory. Every named constant from the environment that appears in this term is called a \textbf{premise} of the declaration. Consider a simple example:

\begin{codebox}
\textit{\color{gray}-- Tactic proof (what the mathematician writes):}\\
\textbf{theorem} \textcolor{red!70!black}{Nat.succ\_pos} (n : Nat) : 0 < n + 1 :=\\
\hspace{1.5em}\textcolor{red!70!black}{Nat.zero\_lt\_succ} n\\[3pt]
\textit{\color{gray}-- The proof term contains:}\\
\textit{\color{gray}--\hspace{0.5em}Nat.zero\_lt\_succ\hspace{1em}(named constant from environment $\to$ premise)}\\
\textit{\color{gray}--\hspace{0.5em}n\hspace{6.3em}(local variable, bound by theorem $\to$ NOT a premise)}
\end{codebox}

\noindent
The distinction is precise: a premise is an \emph{external} declaration---a theorem, definition, constructor, or other named constant already registered in the environment---not a local variable or bound parameter. In the example above, \decl{Nat.zero\_lt\_succ} is a premise because it is a theorem proved elsewhere; \texttt{n} is not, because it is a locally bound variable. In practice, the kernel records additional infrastructure premises (typeclass instances, notation abbreviations) that are invisible in the tactic script but present in the compiled term. It is these mechanically extracted premises that form the edges of $G_{\mathrm{thm}}$.

\begin{definition}[Declaration graph]
\label{def:thm-graph}
The \emph{declaration graph} is a directed graph $G_{\mathrm{thm}} = (\mathcal{D}, E_{\mathrm{thm}})$, where $\mathcal{D}$ is the set of all declarations---theorems, definitions, abbreviations, inductive types, constructors, opaques, quotients, and axioms---in \module{Mathlib}, and
\[
  E_{\mathrm{thm}} = \{(d_1, d_2) \in \mathcal{D} \times \mathcal{D} \mid d_2 \text{ appears as a premise in the proof term of } d_1\}.
\]
Each node carries two attributes: its \emph{kind} (one of eight declaration types) and its \emph{module} (the source file in which it is defined).
\end{definition}

The eight declaration kinds listed above are Lean~4's fundamental building blocks: theorems ($79.1\%$), definitions ($15.8\%$), abbreviations ($2.2\%$), constructors ($1.5\%$), inductive types ($1.2\%$), and three minor kinds (Table~\ref{tab:kind-overview} in Appendix~\ref{app:decl-detail}). We highlight two structural extremes below. The distinct roles these kinds play in the dependency structure are analyzed quantitatively in \S\ref{sec:thm-types}.

\paragraph{Theorems as edge producers; axioms as pure sinks.}
In $G_{\mathrm{thm}}$, a declaration is an \emph{edge producer} if it has high out-degree (its proof invokes many premises) and a \emph{pure sink} if it has zero or near-zero out-degree but is cited by many others (high in-degree). Theorems constitute $79.1\%$ of all declarations and are the primary edge producers. Lean also accepts the keyword \texttt{lemma}, which is stored identically as a \texttt{theorem} in the compiled environment; our dataset does not distinguish the two. A theorem's out-edges point to the premises invoked in its proof, while its in-edges come from later theorems that cite it:

\begin{codebox}
\textbf{theorem} \textcolor{red!70!black}{Nat.add\_comm} (a b : Nat) : a + b = b + a := \ldots
\end{codebox}

\noindent
At the opposite extreme, \module{Mathlib} rests on exactly three axioms---\decl{propext} (propositional extensionality), \decl{Classical.choice} (the axiom of choice), and \decl{Quot.sound}---which have near-zero out-degree but are transitively required by all $308{,}129$ declarations. Among the three, \decl{propext} alone accounts for $23{,}226$ incoming edges.

\paragraph{Abbreviations as invisible hubs.}
Abbreviations (\texttt{abbrev}) are definitionally transparent shorthands that the elaborator always unfolds, acting as notational glue---typically wrapping typeclass projections or operator notation:

\begin{codebox}
\textbf{abbrev} \textcolor{red!70!black}{OfNat.ofNat} (n : Nat) [inst : OfNat a n] : a := inst.ofNat
\end{codebox}

\noindent
The most-cited node in the entire graph, \decl{OfNat.ofNat} (in-degree $89{,}936$), is an abbreviation: every declaration that mentions a numeric literal depends on it. Because the elaborator inserts these references automatically, abbreviations have the most extreme in-degree distribution among non-axiom types---their citation counts reflect the machinery of the proof assistant, not the structure of mathematical reasoning. By contrast, the most-cited \emph{inductive types}---\decl{CategoryTheory.Category} ($44{,}487$), \decl{Real} ($26{,}637$), \decl{Module} ($23{,}404$)---form a layer of mathematical infrastructure whose centrality reflects genuine mathematical importance. The most-cited \emph{constructor}, \decl{Eq.refl} ($69{,}580$), straddles both layers: it is the sole way to introduce an equality proof, making it simultaneously a language primitive and a mathematical necessity.

\paragraph{Edge directionality.}
Not all declaration kinds participate equally in $G_{\mathrm{thm}}$. An edge $(d_1, d_2)$ means that $d_1$'s proof or body \emph{cites} $d_2$; it does not imply an edge in the reverse direction. For example, the theorem \decl{Nat.lt\_irrefl} proves the irreflexivity of strict ordering on natural numbers: no $n$ satisfies $n < n$.

\begin{codebox}
\textit{\color{gray}-- Irreflexivity of < on natural numbers}\\
\textbf{theorem} \textcolor{red!70!black}{Nat.lt\_irrefl} (n : Nat) : Not (n < n) :=\\
\hspace{1.5em}\textcolor{red!70!black}{Nat.not\_succ\_le\_self} n
\end{codebox}

\noindent
The proof invokes \decl{Nat.not\_succ\_le\_self}, which establishes that $n + 1 \le n$ is impossible (since $<$ on \decl{Nat} is defined as $a < b \iff a + 1 \le b$, the statement $n < n$ reduces to $n + 1 \le n$). The kernel also records two infrastructure premises: \decl{LT.lt}, the abbreviation that provides the $<$~notation, and \decl{instLTNat}, the typeclass instance that connects $<$ to \decl{Nat}. This gives three edges:

\begin{figure}[H]
\centering
\begin{tikzpicture}[
  every node/.style={font=\small\ttfamily, inner sep=2pt, text=red!70!black},
  dot/.style={circle, fill=red!70!black, minimum size=3pt, inner sep=0pt},
  dep/.style={->, >=Stealth, thin, red!70!black},
]
  % Premises (cited) at top
  \node[dot, label={[align=center]above:{not\_succ\_le\_self\\[-2pt]{\normalfont\sffamily\scriptsize\color{gray}theorem}}}] (thm) at (-3.2,0) {};
  \node[dot, label={[align=center]above:{LT.lt\\[-2pt]{\normalfont\sffamily\scriptsize\color{gray}abbrev}}}]              (ab)  at (0,0) {};
  \node[dot, label={[align=center]above:{instLTNat\\[-2pt]{\normalfont\sffamily\scriptsize\color{gray}definition}}}]       (def) at (3.2,0) {};
  % Citing theorem at bottom
  \node[dot, label=below:{Nat.lt\_irrefl}] (t) at (0,-1.8) {};
  \draw[dep] (t) -- (thm);
  \draw[dep] (t) -- (ab);
  \draw[dep] (t) -- (def);
\end{tikzpicture}
\caption{Premise edges for the theorem \decl{Nat.lt\_irrefl}: the mathematical premise \decl{not\_succ\_le\_self} (a theorem) and two infrastructure premises \decl{LT.lt} (an abbreviation providing $<$ notation) and \decl{instLTNat} (a typeclass instance). The edges are directed: none of the three premises cites \decl{lt\_irrefl} in return.}
\label{fig:thm-single-edge}
\end{figure}

\noindent
The eight kinds arrange into a clear hierarchy of edge production and consumption (Table~\ref{tab:inter-kind-flow} in Appendix~\ref{app:decl-detail}). Theorems produce $89\%$ of all edges ($7.50$M of $8.44$M), primarily citing definitions and abbreviations---the computational and notational infrastructure on which proofs are built. Axioms and quotients are near-pure sinks (combined out-degree $= 3$). The graph is nearly acyclic; the only bidirectional edges ($4{,}713$ pairs) link inductive types to their own constructors (e.g., \decl{Nat}~$\leftrightarrow$~\decl{Nat.succ}).

\medskip
\noindent
The \decl{Nat.lt\_irrefl} example above showed a single theorem's premises; we now turn to a \emph{pair} of theorems to illustrate how shared premises create the hub structure that dominates $G_{\mathrm{thm}}$. Consider \decl{Nat.add\_left\_comm}, which proves $a + (b + c) = b + (a + c)$ for natural numbers. A tactic-level proof reads:

\begin{codebox}
\textbf{theorem} \textcolor{red!70!black}{Nat.add\_left\_comm} (a b c : Nat) :\\
\hspace{2em}a + (b + c) = b + (a + c) := \textbf{by}\\
\hspace{1.5em}rw [$\leftarrow$ \textcolor{red!70!black}{Nat.add\_assoc}, \textcolor{red!70!black}{Nat.add\_comm} a b, \textcolor{red!70!black}{Nat.add\_assoc}]
\end{codebox}

\noindent
The \texttt{lean-training-data} tool extracts premises not from the tactic script but from the compiled \emph{proof term}.\xinze{LeanScout pipeline: could cross-validate} For this theorem, the extracted premises are:

\smallskip
\begin{center}
\small
\begin{tabular}{ll}
\textit{Mathematical premises:} & \decl{Nat.add\_comm}, \decl{Nat.add\_assoc} \\
\textit{Equality infrastructure:} & \decl{Eq.refl}, \decl{Eq.symm}, \decl{Eq.mpr} \\
\textit{Typeclass instances:} & \decl{instHAdd}, \decl{HAdd.hAdd}, \decl{instAddNat}
\end{tabular}
\end{center}
\smallskip

\noindent
The first two are the mathematical lemmas that the proof logically depends on. The remaining six are \emph{language infrastructure}: equality combinators that the kernel requires to compose rewriting steps, and typeclass instances that elaborate the \texttt{+}~notation into the concrete function \decl{Nat.add}. Every edge in $E_{\mathrm{thm}}$ records one such dependency, so this single theorem contributes eight edges.

A closely related theorem, \decl{Nat.add\_right\_comm} ($a + b + c = a + c + b$), depends on \emph{exactly the same} eight premises. The resulting fragment of $G_{\mathrm{thm}}$ is shown in Figure~\ref{fig:thm-fragment}: both theorems share all three displayed premises, and \decl{Eq.refl}---a constructor with in-degree $69{,}580$---is the single most-cited declaration in the entire graph. This shared-hub pattern is the mechanism by which a small number of foundational declarations accumulate enormous in-degree.

\begin{figure}[H]
\centering
\begin{tikzpicture}[
  every node/.style={font=\small\ttfamily, inner sep=2pt, text=red!70!black},
  dot/.style={circle, fill=red!70!black, minimum size=3pt, inner sep=0pt},
  dep/.style={->, >=Stealth, thin, red!70!black},
]
  % Top row: shared premises (cited)
  \node[dot, label=above:{add\_comm}]  (ac) at (-3,0) {};
  \node[dot, label=above:{add\_assoc}] (aa) at (0,0)  {};
  \node[dot, label=above:{Eq.refl}]    (er) at (3,0)  {};

  % Bottom row: two citing theorems
  \node[dot, label=below:{add\_left\_comm}]  (alc) at (-2.2,-2) {};
  \node[dot, label=below:{add\_right\_comm}] (arc) at (2.2,-2)  {};

  % Edges from add_left_comm
  \draw[dep] (alc) -- (ac);
  \draw[dep] (alc) -- (aa);
  \draw[dep] (alc) -- (er);

  % Edges from add_right_comm
  \draw[dep] (arc) -- (ac);
  \draw[dep] (arc) -- (aa);
  \draw[dep] (arc) -- (er);
\end{tikzpicture}
\caption{A fragment of $G_{\mathrm{thm}}$: the theorems \decl{add\_left\_comm} and \decl{add\_right\_comm} both cite the premises \decl{add\_comm}, \decl{add\_assoc}, and \decl{Eq.refl}. Infrastructure premises (\decl{instHAdd}, \decl{Eq.symm}, etc.) are omitted for clarity. The shared node \decl{Eq.refl} (in-degree $69{,}580$) exemplifies the hub phenomenon analyzed in \S\ref{sec:thm-degree}. All names have the \texttt{Nat.}\ prefix removed except \decl{Eq.refl}.}
\label{fig:thm-fragment}
\end{figure}

\noindent\textbf{Note.}
The module \module{Algebra.Group.Defs} contains dozens of declarations. In $G_{\mathrm{module}}$, this entire module is a single vertex; in $G_{\mathrm{thm}}$, each declaration within it becomes a separate vertex, and each premise reference becomes a separate edge. A single import edge in $G_{\mathrm{module}}$ thus unfolds into potentially hundreds of premise edges in $G_{\mathrm{thm}}$---this is why $|E_{\mathrm{thm}}|$ exceeds $|E_{\mathrm{module}}|$ by a factor of $360$. This expansion is itself the transition from process (the human decision to bundle declarations into a module) to product (each logical dependency made explicit).

At the commit analyzed, $|\mathcal{D}| = 308{,}129$ and $|E_{\mathrm{thm}}| = 8{,}436{,}366$. In comparison with the module-level module graph of \S\ref{sec:import-graph} ($|\mathcal{M}| = 7{,}563$ nodes, $|E_{\mathrm{module}}| = 23{,}570$ edges), the theorem--premise graph has roughly $40\times$ more nodes and $360\times$ more edges. The enormous increase in edge count reflects the fact that a single module import implicitly carries hundreds of individual premise dependencies; $G_{\mathrm{thm}}$ makes each of these dependencies explicit.

The graph is almost entirely connected: $308{,}054$ of $308{,}129$ nodes ($99.98\%$) belong to a single weakly connected component. The remaining $75$ nodes distribute across $71$ tiny components, of which $69$ are singletons. This near-perfect connectivity---stronger even than the module-level graph, which is exactly connected---reflects the deep logical interdependence of mathematical declarations.

\paragraph{Declaration Types as Structural Strata}
\label{sec:thm-types}

The eight declaration kinds occupy distinct structural niches in $G_{\mathrm{thm}}$ (Table~\ref{tab:thm-kind-degree} in Appendix~\ref{app:decl-detail}). The pattern is systematic: axioms sit at the absolute bottom of the dependency chain with the highest average in-degree ($7{,}823$) and near-zero out-degree; theorems occupy the opposite pole with low in-degree (mean~$5.5$) but the highest out-degree (mean~$30.8$, max~$522$). Among the three axioms, \decl{propext} (propositional extensionality) accounts for $23{,}226$ incoming edges---it is invoked in roughly one out of every thirteen declarations. \decl{Classical.choice} ($169$) and \decl{Quot.sound} ($73$) are invoked far less frequently, suggesting that the vast majority of \module{Mathlib}'s reasoning is constructive at the proof-term level, with classical logic concentrated in specific branches.

Inductive types have high in-degree (mean~$273$, median~$31$) and very low out-degree (mean~$4.1$). These are the structural skeletons of mathematics: type definitions that many theorems depend on but that themselves depend on little. Their in-degree top~$10$ reads as a catalogue of core mathematical concepts: \decl{CategoryTheory.Category} ($44{,}487$), \decl{Real} ($26{,}637$), \decl{TopologicalSpace} ($25{,}310$), \decl{CategoryTheory.Functor} ($24{,}295$), \decl{Module} ($23{,}404$), \decl{CommRing} ($20{,}837$).

Theorems occupy the opposite pole: low in-degree (mean~$5.5$, median~$1$) and the highest out-degree (mean~$30.8$, max~$522$). Theorems are \emph{consumers} of dependencies, not providers. The few high-in-degree theorems are not deep mathematical results but equality reasoning lemmas: \decl{Eq.trans} ($58{,}686$), \decl{of\_eq\_true} ($48{,}801$), \decl{congrFun'} ($45{,}516$), \decl{Eq.symm} ($42{,}542$). These are workhorses of Lean's elaborator and tactic framework---proof infrastructure, not mathematical content.

Abbreviations (\texttt{abbrev}) have the most extreme in-degree distribution among non-axiom types: mean~$438$, maximum~$89{,}936$ (\decl{OfNat.ofNat}). These are type-class coercion paths that the elaborator inserts automatically during type-checking. Their high citation counts reflect the machinery of the proof assistant, not the structure of mathematical reasoning.

\noindent\textbf{Two infrastructure layers.}
\label{rem:two-layers}
The degree-by-kind analysis reveals that the hub structure of $G_{\mathrm{thm}}$ decomposes into two distinct infrastructure layers:
\begin{enumerate}
\item \textbf{Language infrastructure.} The highest-cited declarations---\decl{OfNat.ofNat} ($89{,}936$), \decl{Eq.refl} ($69{,}580$), \decl{DFunLike.coe} ($65{,}437$), \decl{Eq.mpr} ($62{,}942$)---are artifacts of Lean's type system: numeral coercions, equality constructors, function-like coercions. Their prominence is determined by the design of the proof assistant, not by the logical structure of mathematics. These nodes are invisible in $G_{\mathrm{module}}$ (which treats entire modules as atoms) and constitute the first layer of process-trace in the product.

\item \textbf{Mathematical infrastructure.} The highest-cited \emph{inductive types}---\decl{CategoryTheory.Category}, \decl{Real}, \decl{TopologicalSpace}, \decl{Module}, \decl{CommRing}---are core mathematical concepts. Their prominence reflects the logical structure of the formalized mathematics itself: category theory, analysis, and algebra are genuinely foundational in the dependency sense. This layer is the product's own structure, free of proof-assistant artifacts.
\end{enumerate}
The module graph of \S\ref{sec:import-graph} conflates these two layers (a module contains both language infrastructure and mathematical content). The theorem--premise graph separates them, enabling the analyst to distinguish what belongs to mathematics from what belongs to its medium of expression.

\paragraph{The Leaf Theorem Phenomenon}
\label{sec:thm-leaves}

Of the $243{,}725$ theorems in $G_{\mathrm{thm}}$, $109{,}088$ ($44.8\%$) have in-degree zero: they are never cited by any other declaration. These \emph{leaf theorems} form the outermost surface of the library---the boundary where formalized knowledge meets the external world. The zero-citation rate varies dramatically across namespaces (Table~\ref{tab:thm-leaves} in Appendix~\ref{app:decl-detail}): namespaces generated by \texttt{@[to\_additive]} reach $80$--$92\%$, while analysis workbench namespaces (\ns{Primrec}, \ns{AnalyticAt}, \ns{HasFDerivWithinAt}) stay below $18\%$---nearly every theorem in these feeds into downstream proofs.

\noindent\textbf{Leaf theorems as API surface.}
\label{rem:leaf-theorems}
The $44.8\%$ zero-citation rate means that nearly half of all theorems in $G_{\mathrm{thm}}$ are leaves. These leaves constitute the \emph{API surface} of \module{Mathlib}---the interface through which the library's knowledge is consumed by external users (mathematicians applying results, AI systems selecting premises, downstream projects building on Mathlib). Their existence is not waste but function: a library that formalized only internally-reused lemmas would be useless to its consumers.

The dichotomy between high-zero-rate namespaces (systematically generated mirrored API) and low-zero-rate namespaces (hand-crafted building blocks) reveals a second process-trace in the product. The \texttt{@[to\_additive]} machinery leaves a distinctive structural fingerprint---entire subtrees of leaf theorems that inflate the graph's surface area without deepening its logical content. This automation artifact is invisible at the module level (a module contains both original and mirrored theorems) but becomes visible at the theorem level, where individual declarations can be classified by their provenance.

\paragraph{Cross-Namespace Flow}
\label{sec:thm-cross-namespace}

We decompose the $8{,}436{,}366$ edges of $G_{\mathrm{thm}}$ into three tiers based on the relationship between the source and target declarations' names. Here ``module'' is the source file containing the declaration, and ``namespace'' is the first component of its fully qualified name (e.g., the namespace of \decl{Nat.lt\_irrefl} is \ns{Nat})---a declaration-name prefix, distinct from the directory-based $\mathrm{dir}(m)$ of \S\ref{sec:cross-namespace}.

Only $9.6\%$ of edges remain within the same module file, and another $9.6\%$ cross module boundaries while staying within the same top-level namespace. The majority---$50.9\%$---cross namespace boundaries entirely (Table~\ref{tab:thm-edge-breakdown} in Appendix~\ref{app:decl-detail}). In $G_{\mathrm{module}}^{-}$, $62.9\%$ of edges stay within the same namespace (\S\ref{sec:cross-namespace}); in $G_{\mathrm{thm}}$, at most $19.2\%$ do. This threefold increase in cross-namespace density arises because module imports package many individual dependencies into a single edge. The module boundary acts as a cognitive container that conceals the true extent of cross-disciplinary dependency.

The strongest cross-namespace coupling is \ns{AlgebraicGeometry} $\leftrightarrow$ \ns{CategoryTheory} ($38{,}042$ edges), followed by \ns{MeasureTheory} $\leftrightarrow$ \ns{Real} ($15{,}668$ edges); a full ranking is given in Table~\ref{tab:thm-cross-ns-pairs} in Appendix~\ref{app:decl-detail}.

